% ============================================================
%  Pre-registro del protocolo experimental - TFG AutoPipe
%  Autor: Ricardo Antonio Borbon Mena
%
%  PROPOSITO
%  ---------
%  Documento fechado y versionado que fija, ANTES de ejecutar los
%  experimentos, las decisiones de protocolo que el anteproyecto
%  (seccion 4.3) no especifica. No modifica ningun criterio de
%  exito ya comprometido: solo cierra sus silencios.
%
%  REGLA DE ORO
%  ------------
%  Ninguna version de este documento se escribe DESPUES de la
%  ejecucion que gobierna:
%    v0.1  semana 1  - decisiones de protocolo (este documento)
%    v0.2  semana 5  - antes de correr el cribado de sensibilidad
%    v1.0  semana 9  - antes de ejecutar el leave-one-dataset-out
%
%  Cada version se cierra con un commit propio y una etiqueta
%  anotada de git (prereg-v0.1, prereg-v0.2, prereg-v1.0).
%
%  Compilar con:  pdflatex Preregistro.tex
% ============================================================

\documentclass[11pt,letterpaper]{article}

\usepackage[utf8]{inputenc}
\usepackage[T1]{fontenc}
\usepackage[spanish,es-tabla]{babel}
\usepackage{lmodern}
\usepackage[letterpaper,top=2.4cm,bottom=2.4cm,left=2.8cm,right=2.5cm]{geometry}
\usepackage{xcolor}
\usepackage{colortbl}
\usepackage{booktabs}
\usepackage{array}
\usepackage{tabularx}
\usepackage{enumitem}
\usepackage{fancyhdr}
\usepackage{titlesec}

\definecolor{tecnavy}{HTML}{003366}
\definecolor{gris}{HTML}{666666}
\definecolor{tablehead}{HTML}{E6EEF5}

\newcolumntype{Y}{>{\raggedright\arraybackslash}X}
\newcolumntype{L}[1]{>{\raggedright\arraybackslash}p{#1}}

\setlength{\parindent}{0pt}
\setlength{\parskip}{5pt}
\renewcommand{\arraystretch}{1.2}
\setlist[itemize]{leftmargin=1.4em,topsep=2pt,itemsep=3pt}
\setlist[enumerate]{leftmargin=1.6em,topsep=2pt,itemsep=3pt}

\titleformat{\section}{\large\bfseries\sffamily\color{tecnavy}}{\thesection}{0.6em}{}
\titleformat{\subsection}{\normalsize\bfseries\sffamily}{\thesubsection}{0.5em}{}

\pagestyle{fancy}
\fancyhf{}
\renewcommand{\headrulewidth}{0.4pt}
\fancyhead[L]{\small\textit{AutoPipe}}
\fancyhead[R]{\small Pre-registro \preregversion}
\fancyfoot[C]{\thepage}

% --- Identidad de la version --------------------------------
\newcommand{\preregversion}{v0.1}
\newcommand{\preregfecha}{4 de agosto de 2026}
\newcommand{\preregtag}{prereg-v0.1}

\begin{document}

% =========== PORTADA ========================================
\begin{center}
{\LARGE\bfseries\color{tecnavy} Pre-registro del protocolo experimental}\\[6pt]
{\large AutoPipe --- Trabajo Final de Graduación}\\[4pt]
{\small Ricardo Antonio Borbón Mena --- Instituto Tecnológico de Costa Rica}\\[14pt]

\begin{tabularx}{0.92\textwidth}{L{3.6cm} Y}
\toprule
\rowcolor{tablehead} \multicolumn{2}{l}{\textbf{Identificación de esta versión}} \\
\midrule
Versión & \preregversion \\
Fecha de cierre & \preregfecha \\
Estado & Decisiones de protocolo fijadas. Pendientes de versiones posteriores: la
versión congelada del espacio de búsqueda (v0.2) y la ejecución del
\textit{leave-one-dataset-out} (v1.0). \\
Etiqueta de git & \texttt{\preregtag}, etiqueta anotada creada al cerrar esta
versión. Su fecha de creación, registrada por git, es el fechado verificable de
este documento. \\
Forma de cita en el informe & «Pre-registro \preregversion, etiqueta
\texttt{\preregtag}, \preregfecha» \\
\bottomrule
\end{tabularx}

\vspace{6pt}
{\small\color{gris} El documento se identifica por su etiqueta y no por el hash del
commit que lo congela: ese hash se calcula sobre el contenido del documento, de modo
que escribirlo dentro de él lo alteraría. El hash correspondiente se obtiene con
\texttt{git rev-parse \preregtag}.\par}
\end{center}

\vspace{10pt}

% =========== 1. PROPOSITO ===================================
\section{Propósito y alcance de este documento}

Este documento fija, con fecha anterior a la ejecución de los experimentos, las
decisiones de protocolo que el anteproyecto no especifica. Su función es que las
reglas de análisis no puedan ajustarse después de conocer los resultados.

\textbf{Este documento no modifica ningún criterio de éxito del anteproyecto.} La
Sección~4.3 del anteproyecto (corpus y presupuesto, métricas primarias, protocolo
de comparación, criterio de éxito de la hipótesis y criterios de éxito de
ingeniería) permanece vigente sin cambios. Lo que aquí se registra son
precisiones sobre puntos en los que ese texto guarda silencio, más decisiones de
diseño internas que el texto no gobierna. La Tabla~\ref{tab:relacion} deslinda
ambas cosas explícitamente.

\begin{table}[htbp]
\centering
\caption{Deslinde entre lo que fija el anteproyecto y lo que precisa este documento.}
\label{tab:relacion}
\small
\begin{tabularx}{\textwidth}{L{4.1cm} L{4.3cm} Y}
\toprule
\rowcolor{tablehead} \textbf{Asunto} & \textbf{Lo que fija el anteproyecto \S4.3} & \textbf{Lo que precisa este documento} \\
\midrule
Modelo base &
«Un clasificador sensible al preprocesamiento (regresión logística o $k$ vecinos
más cercanos de scikit-learn)», fijo en ambas condiciones. &
Se elige $k$ vecinos más cercanos y se congela su configuración completa
(Sección~\ref{sec:modelo}). El texto autoriza expresamente esta opción. \\
\midrule
Espacio de búsqueda &
Nada: \S4.3 no lo especifica. &
Parametrización relativa (fraccional) y espacio canónico versionado
(Sección~\ref{sec:espacio}). \\
\midrule
Corpus &
Criterios de inclusión: licencia abierta, 500--50\,000 instancias, $\leq 100$
características, sin imagen ni texto embebidos, $\leq 10$ clases. Entre 15 y 25
datasets. 50 ensayos por dataset. &
Cómo se eligen los 15--25 entre todos los elegibles: desempate por sensibilidad
al preprocesamiento (Sección~\ref{sec:corpus}). Los criterios de inclusión no se
alteran ni se añaden nuevos. \\
\midrule
Semillas &
«El mismo presupuesto de 30 ensayos y las mismas semillas» en ambas condiciones. &
Cuántas semillas, cuáles literalmente, y cómo se agregan
(Sección~\ref{sec:lodo}). \\
\midrule
Ensayos hasta el objetivo &
«Número de ensayos necesarios para alcanzar al menos el 95\% del mejor valor de
métrica hallado por el arranque en frío con el presupuesto completo». &
Respecto a qué corrida en frío se define el objetivo, si los ensayos inyectados
cuentan, y qué se hace cuando una condición no alcanza el objetivo
(Sección~\ref{sec:lodo}). \\
\midrule
Criterio de éxito &
Reducción $\geq 30\%$ de la mediana de ensayos hasta el objetivo, con
significancia según Wilcoxon pareada por dataset, $\alpha = 0{,}05$. &
Lateralidad de la prueba y análisis secundarios, que no sustituyen al primario
(Sección~\ref{sec:analisis}). \\
\bottomrule
\end{tabularx}
\end{table}

% =========== 2. MODELO BASE =================================
\section{Modelo base congelado}
\label{sec:modelo}

El modelo base es \texttt{KNeighborsClassifier} de scikit-learn, con esta
configuración literal, idéntica en todas las condiciones, datasets y semillas:

\begin{center}
\texttt{KNeighborsClassifier(n\_neighbors=5, weights="uniform",\\
metric="minkowski", p=2, algorithm="brute")}
\end{center}

\textbf{Justificación de la elección.} La Sección~4.3.3 del anteproyecto autoriza
expresamente regresión logística o $k$ vecinos más cercanos. Se elige el segundo
porque opera sobre distancias, de modo que el escalado, la imputación y la
selección de características afectan la métrica de forma sustantiva. Un modelo
base insensible al preprocesamiento produciría un efecto piso: si la métrica
apenas varía entre pipelines, el objetivo del 95\% se alcanza en el primer
ensayo de ambas condiciones y la comparación queda vacía por construcción,
independientemente de la calidad del recomendador. La elección de $k$ vecinos más
cercanos es la mitigación más económica de ese riesgo.

\textbf{Justificación de} \texttt{algorithm="brute"}\textbf{.} Los algoritmos
basados en árboles degradan su ventaja a un centenar de dimensiones y su
construcción puede introducir variación entre ejecuciones. Forzar fuerza bruta
elimina esa fuente de variabilidad y sirve al criterio de reproducibilidad de la
Sección~4.3.5 del anteproyecto (la misma semilla debe producir el mismo pipeline
ganador).

\textbf{Consecuencia sobre el costo.} El modelo invierte el perfil de cómputo:
entrenar es trivial y predecir crece cuadráticamente con el número de instancias,
de modo que el costo total queda dominado por los datasets cercanos al techo de
50\,000 instancias. La medición de ese costo se realiza en la semana 6 y se
registra en la versión v0.2 de este documento.

Ningún hiperparámetro del modelo base entra en el espacio de búsqueda: el
anteproyecto exige un modelo base fijo para que las diferencias observadas sean
atribuibles al pipeline de preprocesamiento.

% =========== 3. ESPACIO DE BUSQUEDA =========================
\section{Espacio de búsqueda}
\label{sec:espacio}

\subsection{Parametrización relativa}

Todo parámetro cuyo dominio dependa del dataset se expresa en términos
\textbf{relativos} y no absolutos. En particular, la selección de
características se parametriza como fracción de columnas conservadas, no como un
número absoluto de columnas.

\textbf{Razón.} El arranque en caliente transfiere configuraciones entre
datasets. Un parámetro absoluto no es transferible: conservar diez columnas
significa algo distinto en un dataset de doce características que en uno de cien,
y el valor podría ser inválido en el dataset receptor. La parametrización
relativa es condición necesaria para que una configuración recomendada sea
ejecutable en un dataset que el recomendador nunca vio.

\subsection{Espacio canónico y versionado}

El espacio de búsqueda tabular es único, idéntico para todos los datasets del
corpus, y se identifica con un número de versión que se registra junto a cada
entrada de la base de meta-conocimiento y junto a cada resultado experimental.

\textbf{Razón.} En la implementación del motor de búsqueda, los parámetros que
Optuna registra por ensayo son \emph{índices enteros} sobre la lista de opciones
de cada paso, más un índice adicional que representa omitir el paso. Un índice no
tiene significado por sí mismo: depende del orden de la lista. Si el espacio
cambiara de claves o de orden entre datasets, una configuración inyectada como
arranque en caliente designaría transformaciones distintas de las previstas, sin
producir ningún error visible. El versionado hace detectable esa condición:
una recomendación proveniente de una versión del espacio distinta de la vigente
se rechaza en lugar de ejecutarse.

La composición concreta del espacio y su identificador de versión se congelan en
la semana 4 y se registran en la versión v0.2 de este documento.

% =========== 4. CORPUS ======================================
\section{Corpus y cribado de sensibilidad}
\label{sec:corpus}

Los criterios de inclusión son exclusivamente los de la Sección~4.3.1 del
anteproyecto. Este documento no añade, elimina ni relaja ninguno. Lo que fija es
el procedimiento para elegir los 15--25 datasets del corpus entre todos los que
resultan elegibles, que el anteproyecto no especifica.

\subsection{Procedimiento, fijado antes de ejecutarse}

\begin{enumerate}
  \item \textbf{Universo.} Todos los datasets que satisfacen los criterios de
        inclusión de la Sección~4.3.1, sin excepción.

  \item \textbf{Muestreo.} Doce pipelines por dataset candidato, extraídos
        uniformemente al azar del espacio de búsqueda canónico de la
        Sección~\ref{sec:espacio}, incluida la opción de omitir cada paso, con
        semilla fija \texttt{4242}.

  \item \textbf{Evaluación.} El modelo base congelado de la
        Sección~\ref{sec:modelo}, la métrica primaria de clasificación del
        anteproyecto (macro-F1), y la misma partición única de entrenamiento y
        validación que emplea el protocolo de búsqueda.

  \item \textbf{Medida de sensibilidad.} El rango interquartílico (IQR) de los
        doce valores de macro-F1. Se registran además, con carácter descriptivo,
        el rango total ($\max - \min$) y el margen sobre el pipeline vacío
        (mejor pipeline aleatorio menos el desempeño sin preprocesamiento).

        \textit{Razón de la elección:} con doce muestras, el rango total lo
        determina un único pipeline degenerado. El IQR mide si \emph{elegir
        entre} pipelines cambia el resultado, que es precisamente lo que
        gobierna la variable «ensayos hasta el objetivo»; el margen sobre el
        pipeline vacío mide si el preprocesamiento \emph{ayuda}, que es una
        pregunta distinta y se reporta aparte.

  \item \textbf{Regla de desempate.} Los candidatos se ordenan por IQR
        descendente y se seleccionan los primeros hasta completar el corpus,
        sujeto a dos restricciones de diversidad que se aplican \emph{después}
        del orden:
        \begin{itemize}
          \item cobertura de los tres regímenes de tamaño (centenas, millares y
                decenas de millares de instancias);
          \item al menos cinco datasets con variables categóricas y al menos
                cinco con valores faltantes, para que la codificación y la
                imputación del espacio de búsqueda queden ejercitadas.
        \end{itemize}
        Cuando una restricción obliga a omitir un candidato mejor ubicado en el
        orden, se registra cuál y por qué.

  \item \textbf{Umbral del subconjunto sensible.} Un dataset se considera
        sensible al preprocesamiento si su IQR es mayor o igual a $0{,}02$ de
        macro-F1 (dos puntos porcentuales). Este umbral queda fijado aquí,
        \emph{antes} de observar la distribución de los valores, y es el mismo
        que define la estratificación de la Sección~\ref{sec:analisis}.

  \item \textbf{Persistencia.} Los doce valores de cada candidato se conservan en
        el repositorio, incluidos los de los candidatos descartados, para que la
        selección sea verificable y no una afirmación.
\end{enumerate}

\subsection{Resultados posibles declarados de antemano}

Si todos los datasets seleccionados superan el umbral, el subconjunto sensible
coincide con el corpus completo y la estratificación de la
Sección~\ref{sec:analisis} queda vacía. En ese caso se reporta como tal. El
umbral no se redefine después de observar la distribución para forzar una
partición del corpus.

% =========== 5. PROTOCOLO LODO ==============================
\section{Protocolo de comparación}
\label{sec:lodo}

El protocolo es el de la Sección~4.3.3 del anteproyecto: validación cruzada por
exclusión de dataset, dos condiciones sobre el dataset excluido con un
presupuesto de 30 ensayos cada una, arranque en frío contra arranque en caliente
con cinco configuraciones inyectadas mediante \texttt{enqueue\_trial}. Lo que
sigue precisa sus puntos no especificados.

\subsection{Semillas}

Cada condición se ejecuta con \textbf{tres semillas}, y son las \textbf{mismas
tres} en ambas condiciones: $\{101,\ 202,\ 303\}$. La construcción de la base de
meta-conocimiento, que el anteproyecto define con un presupuesto único de 50
ensayos por dataset, emplea la semilla $42$. El cribado de sensibilidad emplea la
semilla $4242$.

Repetir con tres semillas responde a que una sola corrida por condición produce
una observación pareada ruidosa: el muestreador TPE con 30 ensayos tiene varianza
apreciable entre semillas, y una conclusión que dependa de la semilla que tocó no
es defendible. Esto no altera el presupuesto de 30 ensayos, ni el emparejamiento
por dataset, ni el número de pares de la prueba estadística.

Si la medición de costo de la semana 6 muestra que el cómputo lo permite, el
número de semillas asciende a cinco. La decisión se registra en la versión v0.2
de este documento, antes de ejecutar el experimento, y en ningún caso se toma
después de observar los resultados.

En la implementación, la semilla gobierna tanto la partición de entrenamiento y
validación como el muestreador de Optuna. Como ambas condiciones comparten la
semilla, dentro de cada semilla la comparación entre frío y caliente se realiza
sobre una partición idéntica.

\subsection{Definición del objetivo}

El objetivo se define por pareja \textbf{(dataset, semilla)}: para el dataset $d$
y la semilla $s$, el objetivo es el 95\% del mejor valor de macro-F1 que
alcanza el arranque en frío con la semilla $s$ sobre el dataset $d$ en su
presupuesto completo de 30 ensayos. Los ensayos hasta el objetivo se calculan por
semilla, y la observación pareada de cada dataset es la \textbf{mediana de las
tres semillas} en cada condición.

\textit{Razón:} la alternativa —definir un objetivo único por dataset a partir de
la agregación de las tres corridas en frío— permitiría que el arranque en frío no
alcanzara su propio objetivo en alguna semilla, reintroduciendo la censura en la
condición de referencia. La definición por pareja preserva el emparejamiento de
semillas que exige el anteproyecto y mantiene la propiedad de la
Sección~\ref{sec:censura}.

\subsection{Los ensayos inyectados cuentan}

Las cinco configuraciones inyectadas como arranque en caliente cuentan dentro de
los ensayos hasta el objetivo. El arranque en caliente dispone por tanto de
veinticinco ensayos de exploración propia dentro del presupuesto de treinta.

\textit{Razón:} no contarlas otorgaría al arranque en caliente una ventaja
puramente contable, ajena a la hipótesis, que se investiga sobre el costo total
de la búsqueda.

\subsection{Casos censurados}
\label{sec:censura}

Cuando una condición no alcanza el objetivo dentro del presupuesto, su valor de
ensayos hasta el objetivo se fija en $31$ (presupuesto más uno). La regla se
aplica de forma simétrica a ambas condiciones y el número de casos censurados se
reporta por condición y a nivel de semilla, no de dataset.

\textbf{La censura es asimétrica por construcción, y conviene enunciarlo con
precisión: el arranque en frío no puede censurarse.} El objetivo es el 95\% del
mejor valor que el propio arranque en frío alcanza en su presupuesto completo; si
ese mejor valor ocurre en el ensayo $t \leq 30$, entonces el frío cruzó el 95\%
en algún ensayo menor o igual que $t$. No se trata de un suceso improbable sino
imposible. Por lo tanto la regla de censura solo puede afectar al arranque en
caliente, y su elección requiere una garantía explícita, que es la siguiente.

\textbf{Chequeo de robustez.} El análisis primario se recalcula íntegro con
castigos más severos para los casos censurados: dos veces y diez veces el
presupuesto ($60$ y $300$). Si la conclusión no cambia con ninguna de las tres
reglas, se declara explícitamente en el informe que la conclusión no depende de
la regla de censura. Si cambia con alguna, se reporta con qué regla y en qué
sentido, sin seleccionar la más favorable.

Con la mediana de tres semillas, la regla de censura solo desplaza la observación
de un dataset cuando al menos dos de sus tres semillas quedan censuradas; el
reporte a nivel de semilla hace visible esa situación.

\subsection{Ausencia de poda efectiva}

El objetivo evalúa una única partición y no reporta valores intermedios, de modo
que el podador por mediana configurado en el motor de búsqueda no interviene.
Los treinta ensayos son treinta evaluaciones completas y ninguno queda sin valor
de métrica, por lo que no se requiere una regla para ensayos podados.

\subsection{Ausencia de fuga en el nivel meta}

La normalización del espacio de meta-características y la métrica de distancia
del recomendador se ajustan \textbf{únicamente con los datasets de entrenamiento
de cada pliegue}, nunca con el dataset excluido. Ajustarlas sobre el corpus
completo contaminaría cada pliegue y inflaría la ventaja del arranque en caliente
de un modo difícil de detectar a posteriori. Es el análogo, en el nivel meta, del
ajuste de las transformaciones con estado usando solo el subconjunto de
entrenamiento.

% =========== 6. PLAN DE ANALISIS ============================
\section{Plan de análisis}
\label{sec:analisis}

\subsection{Análisis primario}

Sobre el \textbf{corpus completo}, se contrastan las medianas de ensayos hasta el
objetivo de ambas condiciones. La hipótesis se considera confirmada si se cumplen
\emph{ambas} condiciones de la Sección~4.3.4 del anteproyecto:

\begin{enumerate}
  \item la mediana de ensayos hasta el objetivo del arranque en caliente es al
        menos un 30\% menor que la del arranque en frío; \emph{y}
  \item la prueba de rangos con signo de Wilcoxon, pareada por dataset, resulta
        significativa con $\alpha = 0{,}05$.
\end{enumerate}

La prueba se realiza \textbf{de una cola}, en la dirección de la hipótesis
(reducción), y se reporta además el valor $p$ de dos colas. La lateralidad se fija
aquí, antes de ver los datos: la hipótesis del anteproyecto es direccional
(«reduce la mediana en al menos un 30\%»), lo que justifica la prueba de una
cola, y con $N$ entre 15 y 25 pares la potencia disponible es limitada. Reportar
ambos valores evita que la elección de lateralidad se convierta en un grado de
libertad oculto.

La variable «desempeño a presupuesto reducido» (mejor valor de métrica tras
quince ensayos), definida en la Sección~4.3.3 del anteproyecto, se reporta como
medida descriptiva. No forma parte del criterio de confirmación, que la
Sección~4.3.4 hace depender exclusivamente de los ensayos hasta el objetivo.

\subsection{Análisis secundarios, declarados como tales}

\begin{enumerate}
  \item \textbf{Estratificado por sensibilidad.} La misma prueba sobre el
        subconjunto de datasets con IQR $\geq 0{,}02$, etiquetada como secundaria
        y acompañada de su $N$. Si el subconjunto queda por debajo de diez
        datasets, se reportan el tamaño del efecto y las medianas sin
        sobreinterpretar el valor $p$.
  \item \textbf{Robustez frente a la regla de censura}, según la
        Sección~\ref{sec:censura}.
  \item \textbf{Tabla por dataset}, con sensibilidad, margen sobre el pipeline
        vacío, ensayos hasta el objetivo por semilla y condición, casos
        censurados y desempeño a quince ensayos. Los datasets planos aparecen en
        su propia fila con su valor de sensibilidad visible, no promediados
        dentro de la mediana global.
\end{enumerate}

Ningún análisis secundario sustituye al primario en la conclusión del proyecto.

\subsection{Análisis exploratorios, no confirmatorios}

Se declaran exploratorios, y por tanto no pueden usarse para sostener la
hipótesis:

\begin{itemize}
  \item el análisis de diversidad de los pipelines ganadores de la base de
        meta-conocimiento (semana 7);
  \item la ablación con un recomendador aleatorio como control (semana 11);
  \item cualquier análisis de causas posterior a la observación de los
        resultados.
\end{itemize}

Su valor es interpretativo y así se presentan en el informe.

\subsection{Resultado negativo}

Conforme a la mitigación del riesgo RI1 del anteproyecto, si el criterio de la
Sección~4.3.4 no se alcanza, el resultado se reporta como tal, sin reformular la
hipótesis ni el umbral. El análisis de causas y la generalización del framework
constituyen resultados defendibles por sí mismos.

% =========== 7. FUERA DE ALCANCE ============================
\section{Qué queda fuera de este documento}

\begin{itemize}
  \item La composición concreta del espacio de búsqueda canónico y su
        identificador de versión: se congelan en la semana 4 y se registran en
        v0.2.
  \item El tope de submuestreo para las meta-características del grupo de
        complejidad: se fija en v0.2, a partir de la medición de tiempos de la
        semana 1 y del tamaño real de los datasets del corpus.
  \item La medición del costo de cómputo y la decisión entre tres y cinco
        semillas: semana 6, registrada en v0.2.
  \item La identidad de los datasets del corpus: resultado de aplicar la
        Sección~\ref{sec:corpus}, registrada en v0.2.
  \item La incorporación o el descarte del agrupamiento como extensión deseable:
        se decide en el punto de control de alcance de la semana 7, con los
        criterios del plan de trabajo, y no afecta el protocolo aquí registrado.
  \item Todo resultado experimental. Este documento registra reglas, no
        hallazgos.
\end{itemize}

% =========== 8. HISTORIAL ===================================
\section{Historial de versiones}

\begin{center}
\begin{tabularx}{\textwidth}{L{1.6cm} L{2.4cm} Y}
\toprule
\rowcolor{tablehead} \textbf{Versión} & \textbf{Fecha} & \textbf{Contenido y motivo} \\
\midrule
v0.1 & \preregfecha &
Versión inicial. Fija el modelo base, la parametrización relativa y el versionado
del espacio de búsqueda, el procedimiento de cribado de sensibilidad, las
semillas, la definición del objetivo, el conteo de los ensayos inyectados, la
regla de censura con su chequeo de robustez y el plan de análisis. Se cierra en
la semana 1 para que las decisiones queden fechadas el día en que se tomaron. \\
\midrule
v0.2 & Semana 5 &
Pendiente. Añadirá la versión congelada del espacio de búsqueda, el tope de
submuestreo de meta-características, el costo de cómputo medido con la decisión
entre tres y cinco semillas, y la identidad del corpus seleccionado. Se cierra
antes de ejecutar el cribado de sensibilidad. \\
\midrule
v1.0 & Semana 9 &
Pendiente. Versión congelada previa a la ejecución del
\textit{leave-one-dataset-out}. Se cierra antes de generar cualquier resultado
experimental. \\
\bottomrule
\end{tabularx}
\end{center}

\end{document}
